\documentclass[letter, 10pt]{article}
\usepackage[latin1]{inputenc}
\usepackage[spanish]{babel}
\usepackage{amsfonts}
\usepackage{amsmath}
\usepackage[dvips]{graphicx}
\usepackage{url}
\usepackage[top=3cm,bottom=3cm,left=3.5cm,right=3.5cm,footskip=1.5cm,headheight=1.5cm,headsep=.5cm,textheight=3cm]{geometry}



\begin{document}
\title{Inteligencia Artificial \\ \begin{Large}Estado del Arte: Problema [Nombre Problema]\end{Large}}
\author{[Nombre autor]}
\date{\today}
\maketitle


%--------------------No borrar esta secci\'on--------------------------------%
\section*{Evaluaci\'on}

\begin{tabular}{ll}
Resumen (5\%): & \underline{\hspace{2cm}} \\
Introducci\'on (5\%):  & \underline{\hspace{2cm}} \\
Definici\'on del Problema (10\%):  & \underline{\hspace{2cm}} \\
Estado del Arte (40\%):  & \underline{\hspace{2cm}} \\
Modelo Matem\'atico (20\%): &  \underline{\hspace{2cm}}\\
Conclusiones (15\%): &  \underline{\hspace{2cm}}\\
Bibliograf\'ia (5\%): & \underline{\hspace{2cm}}\\
 &  \\
\textbf{Nota Final (100\%)}:   & \underline{\hspace{2cm}}
\end{tabular}
%---------------------------------------------------------------------------%
\vspace{2cm}


\begin{abstract}
El Truck and Trailer Routing Problem (TTRP) consiste en la planificacion de rutas de reparto a multiples clientes con demandas distintas, en donde
se tiene una cantidad definida de camiones y cargamentos. Ciertos clientes pueden ser servidos por camiones con cargamento y otros exigen ser servidos
por camiones desacoplados (ya sea por razones geograficas o logisticas).

\end{abstract}

\section{Introducci\'on}
El presente informe tiene el objetivo de explicar el Truck and Trailer Routing Problem (TTRP), empezando por su definicion original y
posteriormente desglosando articulos clave para caracterizar el estado del arte del problema. Para cada articulo se hara un resumen breve,
la razon de su eleccion y su impacto en la forma de encontrar soluciones optimas.

Se definira el problema en relacion a sus variables, restricciones, funcion objetivo y parametros. Usando el articulo inicial como
base para luego detallar los variantes presentados en articulos posteriores, o de ser el caso, explicar como estos difieren a la hora
de abordar la solucion al problema. Al tener muchas aplicaciones reales, existen multiples versiones que agregan variables para
acercarlo mas a la realidad, el informe no pretende desviar mucho del problema original si no que enfocarse en la evolucion de las
heuristicas empleadas.

El estudio de este problema ha sido fundamental en la logistica de repartos masivos, ya que en muchos casos es necesario usar multiples
tipos de vehiculo para satisfacer la demanda. El problema original surge en parte por problemas encontrados en la industria lactea, en
donde los caminos rurales podian ser demasiado estrechos, haciendo dificil el transito de camiones con carga. Tambien se ha utilizado en el
sector urbano, particularmente en partes de Brasil donde la distribucion de bebidas y tabaco ocurrian en zonas poblacionalmente densas.

\section{Definici\'on del Problema}
El primer articulo en definir este problema fue "A tabu search method for the truck and trailer routing problem - I-Ming Chao" (Chao 2002),
el cual consistia en encontrar la ruta mas corta para una cantidad determinada de clientes, camiones y cargamentos.

Se conoce la cantidad de clientes desde el inicio, en donde cada uno tiene una demanda explicita desde el inicio, la cual puede ser distinta
entre si. Los vehiculos empleados para el reparto pueden tomar dos formatos: con carga (vehiculos completos) y sin carga, en donde los clientes
pueden ser satisfechos con alguno de esos tipos. La capacidad de los vehiculos y sus cargas es igual para todos, en donde se conoce la cantidad
dispuesta desde el inicio.

Lo que lo diferencia de un "vehicle routing problem" (VRP) tradicional es que aqui los vehiculos pueden desacoplarse para modificar el
abanico de clientes a los cuales este puede atender, estacionandose en la posicion de un cliente y posteriormente recorriendo una "sub-ruta",
para la cual el vehiculo esta obligado a volver antes de retomar su curso normal.

Se exige que la demanda de todos los clientes sea satisfecha, la cual debe ser atendida por un vehiculo una sola vez, osea que no se puede
repartir la demanda de un cliente entre multiples visitas. Los vehiculos tienen una capacidad definida, por ende se deben atender los clientes
dentro del presupuesto de capacidad inicial, para la demanda de cada uno de estos. Finalmente, si un vehiculo completo se desacopla, solo este
puede volver a recoger su carga inicial.

El objetivo es minimizar la distancia recorrida por la flota completa, la cual considera un punto de partida para todos los vehiculos a la cual
todos deben volver una vez que la demanda total sea satisfecha.

Al ser un problema que se construye a partir de un VRP (el cual es NP-hard) un TTRP tiene por lo menos la misma complejidad, por lo que un
algoritmo exacto que pueda resolver el problema de forma optima es poco probable. El primer articulo propone buscar una solucion factible
primero, definiendo un recorrido completo para cada vehiculo, compuesto por rutas con vehiculo con carga y sub-rutas con el vehiculo
desacoplado.

Las soluciones iniciales se construyen a partir de una "semilla", la cual define a que cliente atiende cada vehiculo primero,
los clientes posteriores se eligen iterativamente en base a los mas lejanos a las semillas anteriores y al deposito (se eligen los mas lejanos
primero). Sin embargo, antes de llegar a las soluciones factibles se implementa una relajacion al problema; se permite que los clientes sean
satisfechos de forma fraccionaria, para la cual se aplica un redondeo una vez que se asignen todos los clientes, reduciendo el costo computacional
del algoritmo. Una vez que se asignan los clientes a las rutas, se crea la ruta para clientes que pueden ser atendidos por vehiculos completos,
posteriormente, se analizan los clientes con rutas desacopladas, para los cuales se crea una sub-ruta para el mas cercano a la ruta principal,
si se pueden atender mas clientes, estos se agregan a la sub-ruta, si la capacidad no da o esta muy lejos, se crea otra sub-ruta.

Luego de construir las soluciones iniciales, se aplica el procedimiento "descent improvement", necesario para transformar nuestras rutas
iniciales en soluciones factibles (efecto de la relajacion aplicada).
El primer paso consiste en mover clientes entre las rutas propuestas, para el cual se define una penalizacion que se calcula a partir de
la demanda, la capacidad y el tipo del vehiculo. Se examinara un solo cliente a la vez, el cual no puede moverse si este se encuentra como
punto de estacionamiento o si la ruta destino no es compatible.
El segundo paso consiste en cambiar a dos clientes de dos rutas distintas, usando un sistema de penalizacion similar al punto anterior, en
donde los cambios mantienen la factibilidad en relacion a la capacidad del camion.
En el tercer paso se analizan las sub-rutas, para las cuales se eliminan las aristas que conectan al inicio de la sub-ruta (osea el inicio y
el final) y se remplazan por aquellas que conectan dos clientes consecutivos de la sub-ruta a un cliente en la ruta principal, en terminos
sencillos, se mantienen los clientes de la sub-ruta pero se conectan de forma distinta a la principal.

Para cada solucion inicial, el procedimiento anterior correra multiples veces hasta que ninguno de los pasos se pueda ejecutar, cuando esto
ocurra, se aplicara una limpieza de tipo "2-opt", una optimizacion conocida para problemas de grafos, la cual consiste en eliminar las rutas
"cruazadas" entre si (comunmente usando en problemas de tipo Traveling Salesperson Problem).

Los pasos anteriores generaran una serie de soluciones factibles pero no necesariamente optimas y es aqui donde entra la parte mas interesante
del articulo, se emplea una heuristica de mejora de tipo "tabu search", para la cual se desarrolla una variante nueva llamada "ojective-based tabu
restriction (OTB)". En terminos muy simples, tabu search consiste en mejorar una solucion inicial al desplazarse entre soluciones vecinas, para la
cual elige el vecino que de el valor mas optimo (incluso si es peor al original), para evitar ciclos esta define una tabla "tabu", que guarda los
caminos recientemente recorridos, permitiendo que este diversifique. Por otro lado, la heuristica nueva define un rango de mejora, en donde, cualquier
valor que se desvie de este rango estara prohibido, se evaluaran los vecinos dentro de un numero de iteraciones definidas y si no encuentra soluciones,
incrementara la desviacion permitida. La combinacion de ambas heuristicas hace que las iteraciones no repitan pasos ciclicamente y que ademas escape
de los maximos locales, ya que los caminos van mejorando gradualmente y se van prohibiendo a la vez, abriendo el paso para alcanzar soluciones
aun mejores. Finalmente, se aplica un criterio de aspiracion (AC) que excluya soluciones en la tabla tabu que son peores a la actual.

\section{Estado del Arte}
Debido al gran impacto que tiene este problema, han surgido multiples articulos abordando TTRP y sus variantes. Uno de los mas interesantes se encuentra
en "A note on the truck and trailer routing problem" (Lin, Yu. Chou 2010), el cual plantea una leve variante denominada "relaxed truck and trailer routing
problem" (RTTRP), el cual evalua la posibilidad de variar la cantidad de vehiculos inciales para asi reducir la distancia recorrida por la flota. El paper
hace un benchmark comparandolo directamente con la propuesta del 2002 y la mejor solucion de ese entonces obtenida el 2009, obteniendo resultados consistentemente
mejores que ambos, con un margen promedio de 5 porciento mejor que la ultima.

Anos mas tarde, aparece "A matheuristic for the truck and trailer routing problem" (Villegas, Prins, Prodhon, Medaglia, Velasco 2013) con el fin de emplear
los ultimos metodos de optimizacion hibrida. Las matheuristicas surgen como la alternativa a los metodos exactos, las heuristicas y las metaheuristicas. Este
acercamiento descompone el problema inicial en multiples sub-espacios, restringiendo el espacio de busqueda a vecinos locales y usando programacion matematica
para hallar una mejor solucion. Se emplea un metodo hybrid GRASP x ILS como columnas en una formulacion de set-partitioning. Este nuevo metodo logro los mejores
resultados en terminos de calidad y tiempo de computo para su epoca, abordando tambien RTTP para flotas con tamano variable.

El ano 2025 aparecio el articulo "A Cooperative Algorithm to Tackle the Truck and Trailer Routing Problem" el cual planteaba resolver el problema usando
algoritmos cooperativos, metodo que consiste en separar un algoritmo en multiples procesos independientes que trabajan en subconjuntos de informacion que
los otros procesos desconocen. Dichos espacios se resuelven por si solos y luego van "evolucionando", generando nuevos espacios de soluciones que sirven
para iterar y llegar a nuevos optimos. En particular, este paper emple el algorimo "CoopA", el cual es competitivo con otros del mismo tipo, por lo que
abre las puertas a una nueva area a estudiar.

Finalmente, uno de los ultimos articulos publicados del tema es "Nuevos optimos para el Problema de Rutas de Camiones y Remolques con Restricciones de Capacidad
usando un algoritmo colaborativo", el cual combina multiples aspectos de los descubrimientos mencionados. Esta propuesta hace uso de CORAL, un algoritmo de busqueda
local centrado en resolver el CTTRP. CORAL consta de dos fases, la primera de contruccion, basado en procedimientos GRAASP y una fase posterior de reparacion basada
en procedimientos Hill-climbing. Esta propuesta obtiene optimos iguales a los de la literatura de forma consistente, pero al mismo tiempo siendo considerablemente
mas eficiente.

Es importante notar que hasta la fecha, el problema siempre se aborda con una metodologia de dos pasos: encontrar una solucion factible rapido y posteriormente
mejorarla con heuristicas especificas, evidente en todos los articulos discutidos hasta ahora.

\section{Modelo Matem\'atico}
En esta seccion se describe el modelo matematico del problema planteado en el articulo original del 2002, desglosando las variables, restricciones, funcion objetivo
y parametros en cada seccion.

\subsection{Parametros}
Los parametros del sistema se conocen desde un principio, primero se define la demanda de cada cliente:

q_i: La demanda del cliente i

En donde cada cliente tiene un tipo, el restringe el tipo de vehiculo que puede atenderlo

ti:
0 vehicle customer (v.c.)
1 truck customer (t.c.)

Un v.c. puede ser atendido por camion solo o por camion + remolque.
Un t.c. solamente puede ser atendido por camion solo.

Luego, tenemos la distancia desde un cliente (i) a otro (j)

c_ij = distancia euclidiana entre i y j

Para la flota, cada camion y remolque tienen capacidades aparte e iguales para cada tipo

Q_1: capacidad de cada camion
Q_2: capacidad de cada remolque

Por lo que la capacidad de un vehiculo completo es Q_1 + Q_2.

La capcidad de la flota tambien se conoce de antemano y es fijo, para problemas de RTTRP esto pasa
a ser una variable pero para el articulo original es una constante.

m_1: numero de camniones
m_2: numero de remolques

En donde se establece que m_1 > m_2. La posicion del deposito se omite aqui por simplicidad.

\subsection{Variables}

De forma preliminar, se separan los clientes por ruta.

x_ij = 1 si el cliente i pertenece a la ruta j, 0 si no

El artículo utiliza explicitamente variables x_ij para su problema de asignacion relajado.

Luego, para cada ruta hay que definir lo siguiente:

1. Cuales son los puntos de acople y desacople
2. Que clientes van en un sub-tour
3. En que orden se visitan
4. Cuantos sub-tours hay

\subsection{Retricciones}

1. Todos los clientes deben ser atendidos:

Cada cliente i debe pertencer a una ruta.

La sumatorio de x_ij para todo j debe ser 1.

El articulo utiliza precisamente esta restriccion en su problema de asignacion.

2. No se puede dividir la demanda.

Esta restriccion va de la mano con la anterior ya que si pertenece a una ruta, se fuerza a que se satisfaga una sola vez.

3. Se respetan las capacidades

Una pure truck route debe cumplir:

La sumatorio de q_i <= Q_1

Lo mismo ocurre con un sub-tour:

La sumatorio de q_i <= Q_1 (con i siendo del sub-tour)

Porque durante el sub-tour está circulando solamente el camión.

4. Tipos de cliente.

Se debe respetar que cada tipo de cliente sea atendido por el vehiculo del tipo correspondiente.

\subsection{Funcion Objetivo}

Aqui el objetivo se define simplemente como:

min distancia total recorrida por la flota

Para las conexiones utilizadas:

(i,j)

Y cada una tiene distancia c_ij, conceptualmente:

mini∑​j∑​cij​xij​

La formulacion concreta del articulo utiliza la distancia total recorrida como objetivo, asumiendo que el costo es proporcional a la distancia.

\section{Conclusiones}
Conclusiones RELEVANTES del estudio realizado. Deber\'ia responder a las preguntas: ?`todas las t\'ecnicas resuelven el mismo problema o hay algunas diferencias?, ?`En qu\'e se parecen o difieren las t\'ecnicas en el contexto del problema?, ?`qu\'e limitaciones tienen?, ?`qu\'e t\'ecnicas o estrategias son las m\'as prometedoras?, ?`existe trabajo futuro por realizar?, ?`qu\'e ideas usted propone como lineamientos para continuar con investigaciones futuras?

\section{Uso LLMs}
Indique ac\'a, para las herramientas de LLMs/IA usadas en su proyecto. Para cada herramienta/uso, indicar:
\begin{enumerate}
\item Herramienta utilizada: ChatGPT, Copilot, Gemini. Incluya versi\'on si es posible: Gratuita, plan.
\item Prop\'osito de uso: Generar ideas iniciales, explicar,  conceptos, redactar borradores, corregir estilo o gram\'atica, programar/depurar c\'odigo.
\item Prompts relevantes: Preguntas clave que hicieron, iteraciones relevantes (si cambiaron la pregunta para mejorar resultados)
\item Nivel de edici\'on: ?`Copiaron y pegaron?, ?`reescribieron?, ?`verificaron con otras fuentes?
\item Validaci\'on: Uso de Libros/papers, apuntes de clase, juicio propio.
\end{enumerate}
%Herramienta: ChatGPT (GPT-5.x)
%Uso: Generación de ideas y revisión de redacción
%Prompt principal: “Explica el algoritmo SMS-EMOA con pseudocódigo formal”
%Intervención: El texto fue adaptado, reorganizado y complementado con material de clase
%Validación: Se contrastó con [fuente X]


%------Esta sección puede ser removida ----------%
\section{Bibliograf\'ia}
Indicando toda la informaci\'on necesaria de acuerdo al tipo de documento revisado. Todas las referencias deben ser citadas en el documento.
%------Remover hasta acá ----------%
\bibliographystyle{plain}
\bibliography{Referencias}

\end{document} 
